% 27 --------------------------------------------------------------------------
\begin{frame}{K-means in One Sentence}
  \begin{tikzpicture}
    \node[orangebox,text width=11.5cm,minimum height=1.2cm]
      {K-means learns a finite set of prototypes by alternating nearest-center assignment and mean updates under a chosen Euclidean geometry.};
  \end{tikzpicture}
  \vspace{.55cm}
  \begin{itemize}\spaceditems
    \item The objective explains the updates;
    \item the geometry explains the failure modes;
    \item initialization and preprocessing shape the result;
    \item evaluation must include stability and domain meaning;
    \item JL projection can reduce dimension while approximately preserving the objective.
  \end{itemize}
\end{frame}


